\section{Slack Admission, Dispatch, and Revocation}
\label{sec:control}
\subsection{Slack is a budget with a scope}
An urgent request has a release time, a completion objective, and a current
estimate of its remaining service demand. The architecture distinguishes
first-token latency, per-step latency, and whole-request completion: satisfying
one does not prove the others. A background admission spends only slack
defined at the same boundary as the controller's latency observations.
Speculative tokens, variable search width, and prefill/decode transitions
require the controller to change its foreground class or invalidate the
corresponding calibration.

For the measured sharing experiment, the objective is a step deadline
$D_f=1.06\,Q_{.99}(T_f^{\mathrm{solo}})$. The controller reports both
$Q_{.99}(T_f^{\mathrm{co}})/Q_{.99}(T_f^{\mathrm{solo}})-1$ and deadline
attainment, the fraction of observed steps no longer than $D_f$. These are
different statistics. Passing a p99 ratio test does not establish that every
step met its deadline; observing a finite sample with no misses does not
establish a deterministic deadline guarantee.

The intended admission problem is to choose background utility subject to
foreground and resource constraints:
\begin{equation}
 \max_{\mathcal A,\mathbf r}\; U(\mathcal A,\mathbf r)
 \quad\text{s.t.}\quad
 \widehat T_f+\widehat C_f(\mathcal A,\mathbf r)+G_f\leq D_f.
\end{equation}
Here $\mathcal A$ is the admitted task set, $\widehat C_f$ includes predicted
interference, and $G_f$ reserves revocation and uncertainty margins. Capacity,
execution-resource, and protection constraints apply independently. Utility
can value completed work, throughput, or application priority. A scalar
FLOP count is appropriate only when comparing work whose output and precision
make that count meaningful.

This formulation does not imply that a useful global optimizer is already
available. The implemented policy uses a measured table for a known workload.
Among numerically qualified cooperative configurations whose calibration
p99 meets the budget, it chooses the highest measured background rate.
Unknown configurations are not admitted by extrapolation. Selection is
frozen before held-out samples are collected. If a selected window fails,
the controller revokes background admission for a later window; it cannot
repair a deadline already missed. The experiment tests this narrow selection
rule, not prediction under arbitrary future workload shifts.

\subsection{The cost model must include the foreground}
A useful cost table records an allocation-only baseline in addition to a
co-run result. Suppose a foreground loses $a$ units of latency when its
worker allocation shrinks, even with no background arithmetic. If a co-run
then loses $a+b$, attributing all $a+b$ to memory contention misdiagnoses the
problem. Conversely, assigning a separate SM subset does not remove shared
cache, DRAM, power, or interconnect effects, so attributing the entire change
to lost SMs is also unjustified.

\subsection{Dispatch and safe revocation}
The lease lifecycle is \emph{prepared}, \emph{eligible}, \emph{running},
\emph{revoking}, \emph{checkpointed}, and finally \emph{retired}. Preparation
verifies tensor views, code versions, resource declarations, and storage
availability. Eligibility requires predecessor completion and a valid grant.
Revocation prevents new grants, requests a safe point from already-running
work, and waits for the specified acknowledgement. Retirement is separate:
it also requires that downstream users release the task's state and outputs.

A revocation acknowledgement means that the advertised restart state is
committed and that the promised resources have been returned. It must not
mean merely that one thread read an urgent flag. A resident worker can only
acknowledge its weaker contract: arithmetic issue has stopped, while its
CTA allocation remains. A finite-slice backend acknowledges after its CUDA
stream confirms completion. The architecture attaches this semantic mode
to telemetry so the policy cannot compare unlike acknowledgements as though
they offered identical scheduling opportunities.

Let $T_{\mathrm{pub}}$, $T_{\mathrm{drain}}$, and $T_{\mathrm{ack}}$ denote
control publication, remaining non-preemptible work, and acknowledgement
cost. A software revocation estimate is
\begin{equation}
 T_{\mathrm{revoke}}\approx
 T_{\mathrm{pub}}+\max_i T_{\mathrm{drain},i}+T_{\mathrm{ack}}.
\end{equation}
It becomes a bound only if each constituent is bounded under the admitted
interference and the control path is guaranteed execution capacity. A quantum
defined as K tiles does not provide that guarantee by itself. Cache misses,
context scheduling, and competing traffic can stretch the time of the same
instruction sequence. This is why the specification retains a hardware
pause requirement even when compiler slicing is available.

\subsection{An IPC implementation of admission}
The process-level prototype coordinates one urgent client and a fixed set
of background clients with a shared condition variable. A background process
acquires the condition lock, verifies that no urgent interval is active, and
marks one slice in flight. It releases the lock, launches its finite slice,
and waits for stream completion. It then clears its in-flight mark under
the same lock and notifies waiters. The foreground sets urgency under that
lock and waits until every in-flight mark is clear before submitting its work.

Using one lock for the check and the in-flight transition closes an important
race. If a worker first checked an event and only later published its active
state, the foreground could observe the old idle state and admit itself just
before that worker launched another slice. The foreground would then have a
false acknowledgement. The condition-variable protocol gives admission a
linearization point. It still depends on trusted workers reporting stream
completion honestly; it is not a security boundary against a malicious CUDA
client.

\subsection{A worked admission and revocation sequence}
Consider an urgent decode client and an adapter-update client whose ready
matrix product has been lowered to output tiles and K chunks. This example
specifies the proposed protocol; it is not an additional experiment. Before
sharing begins, the controller qualifies the foreground geometry, prepares
the background executable, allocates its partial accumulators, and resolves
all input dependencies. It records a code epoch and a lease generation with
the execution cursor. Compilation and first-use initialization occur before
the urgent client depends on a bounded revocation interval.

Suppose the foreground announces an eligible sharing interval. The controller
looks up a qualified allocation and checkpoint policy and subtracts its
revocation reserve from the available slack. A positive remainder permits
one background grant. A nonpositive remainder admits no work, even if tensor
activity is low. The grant identifies the task, generation, allowed program
range, and K extent. The backend issues only that range; the next grant may
change its program count without changing which tiles have already committed.
This makes resource resizing a change in future scheduling rather than an
implicit restart of completed arithmetic.

When urgent work becomes ready, the broker closes admission before waiting
for outstanding slices. The finite-slice executor finishes its current launch,
commits the next cursor, and acknowledges stream completion. A stale completion
from a previous lease generation cannot satisfy this acknowledgement. Only
then does the foreground receive the resources promised by this backend.
If the background is instead a resident cooperative worker, the acknowledgement
has a different meaning: its arithmetic pauses, but its register and shared
memory allocation remains occupied. The controller must have admitted that
retained footprint in advance.

After the urgent interval, the controller validates the foreground class and
lease generation again. An unchanged workload may reuse the qualified policy;
a changed context shape or model phase can require a different entry or solo
execution. Resumption consumes the last committed accumulator and cursor,
not the original inputs alone. Once all output tiles finish, a completion
event releases dependent work. Storage is reclaimed only after those consumers
retire. This sequence joins admission, execution, and lifetime management
without treating a returned host call as proof of device completion.

\subsection{Failure, cancellation, and ownership}
A general OS contract needs a terminal failure state as well as a successful
checkpoint state. The specification distinguishes user cancellation, an
expired admission lease, a worker that stops responding, and a device error.
Cancellation first prevents future grants. It does not retroactively erase
stores or DMA operations already issued. An output is published as successful
only after its completion condition holds; a cancelled partial output remains
private until its in-flight accesses are known to have ended.

The proposed recovery record contains the task identifier, executable epoch,
lease generation, last committed cursor, and the ownership of every durable
partial buffer. Replaying from that record is valid only if the associated
inputs and executable remain available and the task's declared side effects
permit replay. Pure tensor stages make this rule tractable. A kernel that
updates external state or participates in a collective needs a separate
commit protocol; the scheduler cannot manufacture exactly-once effects by
relaunching it after a timeout.

A timeout is evidence of missing acknowledgement, not evidence that memory
is safe to reuse. If a worker process dies while device work may remain,
the controller must obtain a lower-layer completion or teardown guarantee
before recycling its allocations. Until then it quarantines the affected
lease or reports failure to the foreground. A proposed deadline guarantee
must budget this failure behavior separately from ordinary checkpoints.
The current trusted IPC prototype checks its local admission lifecycle; it
does not implement durable crash recovery, hostile-client containment, or
a bounded device-reset service. Those requirements define how the architecture
would extend beyond the tested cooperative setting.

\subsection{Progress, fairness, and overload}
Foreground protection can starve background work. A controller that pauses
backgrounds for every full decode step may be safe for the foreground but
economically useless when those steps occupy almost the entire timeline.
The architecture therefore reports background progress alongside deadline
attainment. A long interval of zero progress should trigger reconsideration
of the lease, migration to another device if available, or explicit admission
rejection rather than indefinite retention of scarce memory.

\section{Memory and Resource Governance}
\label{sec:memory}
\subsection{Storage lifetime is independent of execution lifetime}
Memory capacity often determines whether another workload can be admitted
even when arithmetic is available. A process may hold model weights,
request-private KV, shared prefixes, scratch, and checkpoints with very
different lifetimes. The proposed allocator names these classes separately.
It tracks ownership, consumers, mutability, residency, and reclaim cost,
rather than treating all allocated bytes as equally revocable.

An immutable weight image can be retained while its process has no runnable
task. A private KV page belongs to a logical request, not to the current
position of that request in a compacted batch. A shared prefix requires
reference-counted ownership and copy-on-write before divergence. Scratch can
be recycled only after the last producer and consumer complete. A partial
reduction is durable task state until the task retires, even though the
kernel that produced it has returned. These distinctions are essential if
execution is to be paused or reorganized without corrupting storage.

\subsection{Capacity leases and eviction}
The architectural capacity lease includes retained bytes, an eviction class,
and an expected restoration path. Eviction can mean dropping recomputable
scratch, writing checkpoints to host memory, preserving an immutable host
weight image, or refusing new requests whose KV would exceed the pool.
These choices have different costs and correctness requirements. A byte of
unused reserved capacity can be cheap to lend; a byte holding the only copy
of an active request's state may be expensive to reclaim.

Revocation of memory is a transaction. The controller first prevents new
users, then waits for existing device accesses, produces a durable copy if
required, updates the residency map, and finally returns the physical
allocation. An application-visible pointer cannot remain valid after its
backing storage has been reassigned unless the lower layer supplies the
corresponding address-translation guarantee. The design therefore separates
logical handles from physical mappings and requires versioned residency
events for asynchronous consumers.

There is no assumption that paging is always preferable to rejection. Moving
state consumes copy engines and host links, may disturb the foreground's
memory traffic, and can make subsequent restoration exceed its latency
objective. A controller should compare those costs with the remaining useful
background work. It should also account for the cost of re-prefill if KV
state is discarded. Capacity utilization without restoration cost is an
incomplete admission metric.

\subsection{Bandwidth pacing and cache pressure}
The implemented resident worker can impose a minimum interval between CTA
slice starts. It does not accumulate unused credits into a later burst. This
provides a simple logical issue-rate control and prevents a long idle period
from authorizing an unbounded foreground-interfering burst. The descriptor
records operand and checkpoint bytes, but the control does not enforce a
physical DRAM or L2 bandwidth ceiling.

A stronger governor would coordinate budgets across clients and resource
domains. Per-worker pacing alone scales the aggregate allowed traffic with
the number of workers. Reusing the same per-worker quota after doubling the
grid can therefore double the intended issue envelope. The lease specifies
whether a budget is per worker, per process, or per device. Global budgets
require either a shared credit mechanism or hardware support; they cannot
be inferred from independent local sleeps.

\subsection{Copies, links, and power}
Host transfers, peer transfers, collectives, and compute can contend for
shared queues or memory bandwidth even when they use different engines.
The proposed contract tags a transfer with its direction, source and
destination, byte count, deadline class, and whether it is necessary to
unblock urgent work. Copy admission must reserve the corresponding memory
traffic as well as the copy engine. Otherwise a nominally independent DMA
operation can exceed the foreground's memory budget.

Power and frequency are shared state rather than per-task constants. A
background workload can alter device clocks and thereby change both
foreground and background service rates. The experiments use unlocked
clocks and record that limitation; they do not derive a power-aware price
model from isolated timing ratios. The proposed governor would need
calibrated telemetry and an explicit policy for thermal or board-power
constraints. It cannot obtain a per-process power guarantee by adjusting
only the number of submitted CTAs.

\subsection{Distributed extensions and tail composition}
Although the current scope is one 32-GB GPU, the architecture admits a
distributed coordinator above device-local leases. A collective phase must
not be admitted independently on ranks whose progress depends on one another.
Pausing a rank inside an uncheckpointable collective can strand peers and
invalidate a local slack estimate. Collective entry and completion therefore
become explicit gang-scheduling boundaries until a lower layer provides a
stronger protocol.

For illustration, if independent rank completion times have CDF $F(t)$,
the latest of $R$ ranks has CDF $F(t)^R$. A local quantile constraint is not
automatically a job-level constraint. Real ranks are correlated by shared
inputs, topology, and collective timing, so the independence expression is
a model rather than a deployment measurement. The design requires the
coordinator to price the observed job tail and to distinguish actual rank
measurements from projected maxima. No multi-GPU result is claimed in this
paper's single-device experiments.
